\section{Contract Framework}\label{sec:contracts}
% Budget: 2.5--3 pages. Mirrors lean/ContractMonitors/{Basic,Contract,
% Compositional,Interface}.lean; keep both in sync.

We employ the \textsc{Focus} notation~\cite{BroyStolen2001}. A system
$\sys{s} \in \SYSTEM$ interacts with its environment exclusively via typed
channels $c \in \CHANNEL$. $\SYSTEM$, $\CHANNEL$, and $\MESSAGE$ denote the
sets of all systems, channels, and messages, respectively.

\paragraph{Streams and Interfaces.}
Time is divided into discrete intervals (ticks). A \emph{timed stream} over
$M \subseteq \MESSAGE$ is a function $s\colon \natp \to M \cup \{\nomsg\}$,
where $s(t) = \nomsg$ denotes that no message is transmitted in
interval~$t$. A \emph{channel history} for a set of channels
$C \subseteq \CHANNEL$ assigns a timed stream to every channel in~$C$; the
set of all channel histories for $C$ is denoted by $\chan{C}$. For
$x \in \chan{C}$ and $t \in \nat$, we write $\pre{x}{t}$ for the
\emph{prefix} of $x$ comprising the first $t$ intervals, and
$\proj{x}{C'}$ for the \emph{projection} of $x$ to the channels
$C' \subseteq C$.

The \emph{syntactic interface} of a system is a pair $\iface{I}{O}$ of
disjoint sets of input channels $I$ and output channels $O$. Its
\emph{semantic interface} is a function
$f\colon \chan{I} \to \wp(\chan{O})$ that maps every input history to the
set of possible output histories.

\paragraph{Contracts.}
A contract separates the obligations of a system from those of its
environment.

\begin{definition}[Assume/Guarantee Contract]\label{def:contract}
  A \emph{contract} for the syntactic interface $\iface{I}{O}$ is a pair
  $C = \ctr{A}{G}$ of predicates, where the \emph{assumption}
  $A \subseteq \chan{I}$ constrains the input histories and the
  \emph{guarantee} $G \subseteq \chan{I \cup O}$ relates input and output
  histories.
\end{definition}

\begin{definition}[Satisfaction]\label{def:satisfaction}
  A system $\sys{s}$ with syntactic interface $\iface{I}{O}$ and semantic
  interface $f$ \emph{satisfies} the contract $C = \ctr{A}{G}$, written
  $\sys{s} \models C$, if
  $\forall x \in \chan{I}, y \in f(x)\colon A(x) \implies G(x, y)$.
\end{definition}

\noindent
Hence, a contract corresponds to the interface assertion $A \implies G$;
keeping both parts apart is what later allows a violation to be attributed
either to the system or to its environment (cf.\
Sect.~\ref{sec:compositional}).

\paragraph{Refinement and Composition.}
Refinement relates the contracts of adjacent levels of the architecture,
e.\,g., the contract of a system and the contracts of its components.

\begin{definition}[Refinement]\label{def:refinement}
  A contract $C' = \ctr{A'}{G'}$ \emph{refines} a contract
  $C = \ctr{A}{G}$, written $C' \refines C$, if (i)~$A \implies A'$ and
  (ii)~$A \land G' \implies G$.
\end{definition}

\noindent
A refined contract thus assumes less about the environment and, in every
environment admitted by the abstract contract, guarantees at least as much.

% OPEN: the assumption of a composed contract refers to G and hence to output
% channels, whereas Def.~\ref{def:contract} restricts assumptions to inputs.
% Either state Def.~\ref{def:contract} for atomic contracts only, or relax it.
\begin{definition}[Composition]\label{def:composition}
  Let $C_1 = \ctr{A_1}{G_1}$ and $C_2 = \ctr{A_2}{G_2}$ be contracts. Their
  \emph{composition} is the contract $C_1 \comp C_2 = \ctr{A}{G}$ with
  $G \equiv (A_1 \implies G_1) \land (A_2 \implies G_2)$ and
  $A \equiv (A_1 \land A_2) \lor \lnot G$.
\end{definition}

\begin{theorem}[Soundness of Refinement]\label{thm:refinement}
  Refinement is reflexive and transitive. Moreover, if $C' \refines C$ and
  $\sys{s} \models C'$, then $\sys{s} \models C$.
\end{theorem}

\begin{theorem}[Independent Implementability]\label{thm:composition}
  If $\sys{s}_1 \models C_1$ and $\sys{s}_2 \models C_2$, then the
  composition of $\sys{s}_1$ and $\sys{s}_2$ satisfies $C_1 \comp C_2$.
  Furthermore, composition is monotonic with respect to refinement: if
  $C_1' \refines C_1$ and $C_2' \refines C_2$, then
  $C_1' \comp C_2' \refines C_1 \comp C_2$.
\end{theorem}

\noindent
Both theorems are mechanised (cf.\ Sect.~\ref{sec:mechanisation}).
Theorem~\ref{thm:composition} permits components to be developed, replaced,
and---as we show in Sect.~\ref{sec:compositional}---monitored independently
of each other. None of the definitions refers to a particular level of the
architecture; they apply alike to systems of systems, systems, and
components.

\paragraph{Monitorable Contracts.}
A runtime monitor observes a finite prefix of a channel history only.
Consequently, it can establish violations of those properties whose
violation is witnessed by a finite prefix.

\begin{definition}[Monitorable Contract]\label{def:monitorable}
  A \emph{prefix predicate} $p$ over a set of channels $C$ is a decidable
  predicate on the finite prefixes of channel histories in $\chan{C}$. It
  induces the property
  $\always{p} = \{x \in \chan{C} \mid \forall t \in \nat\colon
  p(\pre{x}{t})\}$.
  A contract $\ctr{A}{G}$ for $\iface{I}{O}$ is \emph{monitorable} if
  $A = \always{a}$ and $G = \always{g}$ for prefix predicates $a$ over $I$
  and $g$ over $I \cup O$.
\end{definition}

\noindent
A monitorable contract fixes which channels a monitor has to observe: the
input channels for the assumption, and the channels of the syntactic
interface for the guarantee. The past-time language of
Sect.~\ref{sec:synthesis} provides a syntax for prefix predicates.

\subsection{Contract Patterns for Perception}\label{sec:patterns}
\todo{Patterns as instances of Def.~\ref{def:monitorable}: operational
design domain assumptions, plausibility (range, threshold, cross-data
consistency), temporal consistency, confidence; one specification box per
pattern, instantiated for the running example.}
